\subsection{THETA3-C36: accessible exceptional Weil targets and cycles in every even dimension}\label{C36}

The concrete fourfold of C25 and its cycle from C32 give an
additional family in every even dimension. These targets are
products. The conclusion below does not assert simplicity or
N\'eron--Severi rank one for the particular higher-genus
Jacobians in C24.

\paragraph{The targets and their access from one ball cover.}
Let $J=J(H)$ be the simple fourfold of C25, with
$\operatorname{End}^0(J)=K=\Q(i)$ and
$\operatorname{NS}(J)_\Q=\Q\Theta$. Let
$C_J\in\CH^2(J)$ be the integral cycle constructed in
C32, so that its class lies in the $K$-Weil space and
\[
                 \int_J\operatorname{cl}(C_J)^2=73728.
 \tag{C36a}
\]
For $n\ge2$ put $m=n-2$ and choose $m$ pairwise
nonisogenous non-CM elliptic curves $F_1,\ldots,F_m$
from C16. Define
\[
                  A_n=J\times\prod_{j=1}^m F_j^2.
 \tag{C36b}
\]
On each $F_j^2$ let $i$ act by
$\rho_j(x,y)=(-y,x)$, and use the given action on $J$.
The product polarization has Rosati involution
$i\mapsto-i$. On the two-dimensional holomorphic tangent
space of $F_j^2$, the eigenvalues are $i,-i$, each once.
Thus $A_n$ has dimension $2n$ and $K$-signature $(n,n)$.
It is non-CM because it has the non-CM factor $J$.

There is a connected finite \'etale arithmetic ball cover
$X_n\to Y$ of the same Schoen surface with a surjection
\[
                    \operatorname{Alb}(X_n)
                               \longrightarrow A_n.
 \tag{C36c}
\]
Here the two copies of every $F_j$ are obtained from the
two different curve projections, rather than counted
twice from one Albanese map.

\textbf{Proof of (C36c).}
C24's even-level transfer, applied in C25, supplies a
connected finite \'etale cover $D_0\to C_0$ with a
nonconstant map $D_0\to H$. The C16 construction gives
connected finite \'etale covers $D_j\to C_0$ with
nonconstant maps $D_j\to F_j$. Let $D$ be any connected
component of
$D_0\times_{C_0}\cdots\times_{C_0}D_m$.
It is finite \'etale over $C_0$. Its projection onto every
connected $D_j$ is finite \'etale, has nonempty image,
and that image is open and closed. Hence each projection
is surjective. In particular $D$ maps nonconstantly to
$H$ and to every $F_j$.

The induced homomorphism
\[
                  J(D)\longrightarrow
                         J\times\prod_{j=1}^mF_j
 \tag{C36d}
\]
is surjective. Each component is surjective by pullback
of holomorphic differentials along the corresponding
nonconstant curve map. Its target factors are simple
and pairwise nonisogenous: $J$ has dimension four and
each $F_j$ has dimension one, and the latter were chosen
pairwise nonisogenous. Complete reducibility therefore
shows that their simple rational first-Hodge-structure
images in $H^1(D,\Q)$ have direct sum. Equivalently,
the dual of (C36d) is injective up to isogeny, and the
image of (C36d) has full dimension. An image abelian
subvariety of full dimension is the whole target.

C12 supplies a generically finite map
$Y\to C_0^2$ of degree $125$. Take a connected component
\[
                   X_n\subset Y\times_{C_0^2}(D\times D).
 \tag{C36e}
\]
It is finite \'etale over $Y$, and surjects onto $Y$
by the same open-and-closed argument. Its proper image
in $D^2$ maps surjectively onto $C_0^2$. Since
$D^2\to C_0^2$ is finite, that image has dimension two;
as $D^2$ is irreducible, the image is all of $D^2$.
Thus $X_n\to D^2$ is dominant and generically finite.
Pullback on $H^1$ is injective: for a generically finite
proper map of smooth projective surfaces, its push-pull
on rational cohomology is multiplication by the nonzero
generic degree. The corresponding map
$\operatorname{Alb}(X_n)\to J(D)^2$ is consequently
surjective. Compose its first $J(D)$ factor with all
components of (C36d), and its second factor with the
$F_j$ components. The result is (C36c).

If $e_j=\deg(D_j/C_0)$ and $e=\deg(D/C_0)$, this
construction retains the finite bounds
\[
 e\le\prod_{j=0}^m e_j,\qquad
 \deg(X_n/Y)\le e^2,\qquad
 \deg(X_n/D^2)\le125.
 \tag{C36f}
\]
For the specific sources above one may take
$e_0\le43264$ and $e_j\le4p_j^2$, where $p_j$ is the
odd torsion prime used to construct $F_j$ in C16.
The cover is compact arithmetic because it is finite
\'etale over the compact arithmetic ball surface $Y$.
All maps and varieties are defined over $\overline\Q$;
a finite extension defines the finitely many choices.
\hfill$\square$

\paragraph{Two explicit Weil divisors on an elliptic square.}
Fix $F=F_j$ and write $\omega_1,\omega_2$ for the
pullbacks of $[0]\in H^2(F,\Q)$. Choose a symplectic
integral basis $u,v\in H^1(F,\Q)$ with
$\int_Fu\wedge v=1$, and use subscripts for the two
factors. For the indicated $i$-action,
$u_1+i u_2$ and $v_1+i v_2$ have eigenvalue $i$.
Put
\[
 \begin{split}
 \beta_+&=(u_1+i u_2)\wedge(v_1+i v_2),\\
 \beta_-&=(u_1-i u_2)\wedge(v_1-i v_2).
 \end{split}
 \tag{C36g}
\]
They span the two embedding lines of
$\bigwedge_K^2H^1(F^2,\Q)$. Both have Hodge type
$(1,1)$. More explicitly,
\[
 \beta_\pm=D\pm iE,\qquad
 D=\omega_1-\omega_2,\qquad
 E=\omega_1+\omega_2-[\Delta_F].
 \tag{C36h}
\]
Indeed the difference map $(x,y)\mapsto x-y$ pulls
$u\wedge v$ back to
$(u_1-u_2)\wedge(v_1-v_2)$ and its zero divisor is
the diagonal. This gives exactly the cross term in
(C36h). Thus $D$ and $E$ are integral divisor classes.
Direct exterior multiplication gives
\[
 \beta_+^2=\beta_-^2=0,\qquad
 \int_{F^2}\beta_+\beta_-=-4,\qquad
 D=\tfrac12(\beta_++\beta_-),\qquad
 \int_{F^2}\beta_\pm D=-2.
 \tag{C36i}
\]
These signs will fix both of the quantitative identities
below. Write $D_j$ also for the actual divisor
$\operatorname{pr}_1^*[0]-\operatorname{pr}_2^*[0]$
on the $j$th elliptic square.

\paragraph{An algebraic projector and the resulting rational cycle.}
Use the C13 projector onto the two $K$-Weil lines in
$H^{2n}(A_n,\Q)$. It can be specified without any
choice of an abstract Hodge projector. Put $b=2+i$,
\[
 f_n(t)=(t-b^{2n})(t-\bar b^{2n}),\qquad
 g_n(t)=\prod_{r=1}^{2n-1}
                       (t-b^r\bar b^{2n-r}).
 \tag{C36j}
\]
Both polynomials lie in $\Z[t]$. Their roots are
distinct because $b/\bar b=(3+4i)/5$ is not a root
of unity in $\Q(i)$. Choose rational polynomials
$r_n,s_n$ by the Sylvester-matrix Bezout identity, so
$r_ng_n+s_nf_n=1$ and their denominators divide
$|\operatorname{Res}(f_n,g_n)|$, and put
$e_n=r_ng_n$. Then $e_n([b]^*)$ is the required
cohomological projector, realized by a rational
combination of algebraic graph correspondences.
Define
\[
 Z_n=e_n([b]^*)\left(C_J\times
                              \prod_{j=1}^mD_j\right)
                  \in\CH^n(A_n)_\Q.
 \tag{C36k}
\]
No Chow-idempotence assertion is needed: the polynomial
is idempotent on the specified cohomology, which is
what is used below.

Let $u_+,u_-$ be the two $K$-embedding components of
$\operatorname{cl}(C_J)$, so
$\operatorname{cl}(C_J)=u_++u_-$.
The determinant of the direct sum of the $K$-modules
in (C36b) is the tensor product of their determinants.
The all-$i$ and all-$-i$ terms are therefore exactly
the terms selected by $e_n$. Equations (C36h)--(C36i)
give the explicit class
\[
 \operatorname{cl}(Z_n)=2^{-m}
 \left(u_+\otimes\bigotimes_{j=1}^m\beta_{j,+}
       +u_-\otimes\bigotimes_{j=1}^m\beta_{j,-}\right).
 \tag{C36l}
\]
It is a nonzero pure $K$-Weil class. If
$\pi:A_n\to J$ is projection, then
\[
 \operatorname{cl}\left(
   \pi_*\left(Z_n\cdot\prod_{j=1}^m
                         \operatorname{pr}_j^*D_j\right)\right)
          =(-1)^m\operatorname{cl}(C_J),
 \qquad
 \int_{A_n}\operatorname{cl}(Z_n)^2
          =(-1)^m73728.
 \tag{C36m}
\]
The first identity uses $2^{-m}(-2)^m=(-1)^m$.
For the second, the squares of each single embedding
line vanish, and the two cross terms in (C36l) give
$2^{-2m}(-4)^m\int_J\operatorname{cl}(C_J)^2$.
Every factor has even cohomological degree, so no
additional K\"unneth sign occurs.

The coefficients of $Z_n$ are rational as stated.
One explicit denominator bound is
\[
                 d_n=|\operatorname{Res}(f_n,g_n)|>0.
\]
The Sylvester-matrix adjugate gives the Bezout identity
with coefficients in $d_n^{-1}\Z$, so $d_ne_n$ has
integral coefficients. Thus $d_nZ_n$ is represented
by an integral signed cycle. Its self-intersection
is $d_n^2(-1)^m73728$, and the first identity in
(C36m) acquires the factor $d_n$. Alternatively any
smaller actual denominator of the chosen projector
may be retained explicitly.

\paragraph{Exceptionality of the entire Weil space.}
Let $B=\prod_jF_j^2$. There is no nonzero homomorphism
between $J$ and $B$: simplicity of $J$ and the distinct
simple isogeny types exclude it. The normalized
Poincar\'e description of line bundles on a product
therefore gives
\[
             \operatorname{NS}(J\times B)_\Q
              =\operatorname{NS}(J)_\Q
                         \oplus\operatorname{NS}(B)_\Q.
 \tag{C36n}
\]
Consequently the subalgebra generated by divisor classes
is the tensor product of the two divisor subalgebras.
Multiplication by $\prod_jD_j$ and proper pushforward
along $\pi$ sends it into the divisor algebra of $J$:
expand a polynomial in the two pullback families and
integrate its $B$-part.

On the two-dimensional $K$-Weil space of $A_n$, that
same contraction is an isomorphism onto the Weil space
of $J$. In the two embedding lines its multipliers
are both $(-2)^m$, by (C36i). They are nonzero; the
formula also respects rational conjugation.
C25 proves that the Weil space of $J$ meets its divisor
algebra only in zero. Applying the contraction shows
\[
 W_K(A_n)\ \cap\
       \Q[\hbox{rational divisor classes on }A_n]=0.
 \tag{C36o}
\]
In particular (C36k) is an actual exceptional algebraic
cycle in every dimension $2n$, $n\ge2$.
Applying $[2+i]^*$ gives a second independent rational
Weil class, since $(2+i)^{2n}$ is nonreal. Thus these
two algebraic cycles span the full Weil space.

The construction is an explicit product consequence of
the checked fourfold cycle and the accessible elliptic
factors. It does not prove that the original genus-$2n$
Jacobians of C24 have N\'eron--Severi rank one, and it
uses no general assertion about Weil classes on arbitrary
non-CM abelian varieties. The common ball cover proves
the source realization (C36c); the exceptional cycle
itself is the algebraic correspondence (C36k), rather
than the image of a surface factoring through a product
of two curves.
